\documentclass[
    %draft, % Mit % kommentieren, um Bilder sichtbar zu machen und Links zu aktivieren
    pdftex,
    a4paper,
    oneside,
    parskip,
    numbers=noenddot,
    listof=totoc,
    bibliography=totoc,
    hyperfootnotes=false
]{scrreprt}
\setuptoc{toc}{totoc}

\newcommand{\thesistitle}{Untersuchung der allgemeinen Anwendbarkeit von Fluxkondensatoren}
\newcommand{\thesistype}{B A C H E L O R A R B E I T}
\newcommand{\thesistypedesc}{im Fachbereich Elektrotechnik/Informatik \\
    der Universität Kassel}
\newcommand{\thesisauthorname}{Niklas Hahn}
\newcommand{\thesisauthorhomestreet}{Mittelstraße 21}
\newcommand{\thesisauthorhometown}{34277 Fuldabrück}
\newcommand{\thesisauthormatrikelnumber}{35522661}
\newcommand{\thesisauthoremail}{niklas.hahn@student.uni-kassel.de}
\newcommand{\thesisdepartment}{Fachgebiet Software Engineering}
\newcommand{\thesisfirstreviewer}{Dr.\ Jens Kosiol}
\newcommand{\thesissecondreviewer}{Prof.\ Dr.\ Gerd Stumme}
\newcommand{\thesissupervisor}{Sebastian Copei\ M.Sc.}
\newcommand{\thesisdate}{\today}

% Select input encodung, usually utf8 is the best choice, on windows, \usepackage[latin1]{inputenc} maybe required
\usepackage[utf8]{inputenc}
\usepackage[T1]{fontenc}
\usepackage[ngerman]{babel}
\usepackage{csquotes}
\usepackage{xcolor}
\usepackage[inkscapeformat=png]{svg}

\MakeOuterQuote{"} % Damit ist es möglich, " " zu verwenden ohne Umlaut zu erzeugen
\defaulthyphenchar=127 % Dadurch werden auch Wörter mit Bindestrich getrennt, die schon Bindestriche enthalten.

% geometry
\usepackage[bindingoffset=1cm, left=2.5cm, right=2.5cm, top=2.5cm, bottom=2.5cm]{geometry}

% Headline
\usepackage{fancyhdr}
\pagestyle{fancy}
\renewcommand{\chaptermark}[1]{\markboth{\thechapter\ #1}{}}
\lhead{\leftmark} \rhead{\thepage}
\cfoot{}
\fancypagestyle{plain}{}

\RedeclareSectionCommand[beforeskip=1.5cm,afterskip=1cm]{chapter}

\usepackage[capposition=bottom]{floatrow}
\usepackage{float}

% Colors
\usepackage{color}
\usepackage{colortbl}

% Tables
\usepackage{tabularx}
\usepackage{multirow}
\usepackage{makecell}
\setlength{\tabcolsep}{4pt}

% Drawing graphs etc.
\usepackage{pgf}
\usepackage{tikz}
\usetikzlibrary{arrows,automata}

%add command for circle drawing in prepartion analysis table
\newcommand*\emptycirc[1][1ex]{\tikz\draw (0,0) circle (#1);}
\newcommand*\halfcirc[1][1ex]{%
  \begin{tikzpicture}
    \draw[fill] (0,0)-- (90:#1) arc (90:270:#1) -- cycle ;
    \draw (0,0) circle (#1);
  \end{tikzpicture}}
\newcommand*\fullcirc[1][1ex]{\tikz\fill (0,0) circle (#1);}

% Footnotes
\usepackage{footmisc}

\usepackage{xspace}
\newcommand{\sic}{[\acs{sic}]\xspace}

% math
\usepackage{amsmath}
\usepackage{siunitx}

% lists
\usepackage{paralist}
\usepackage{enumitem}

% Figures
\usepackage{graphicx, wrapfig}

% Hyperlinks
\usepackage[hyphens]{url}
\usepackage{hyperref}
\hypersetup{colorlinks, citecolor=black, linkcolor=black, urlcolor=black}

% Minted
\usepackage[chapter]{minted}
%\usemintedstyle{xcode}
\setminted{frame=single,tabsize=2,linenos,autogobble}

\newmintinline[code]{text}{breaklines}

\newminted[mdcodeblock]{md}{autogobble,frame=none,linenos=false,breaklines}
\usemintedstyle[js]{abap}
\renewcommand\theFancyVerbLine{\normalsize\arabic{FancyVerbLine}}

% list of abbreviations
\usepackage[printonlyused]{acronym}

% Set line pitch
\usepackage{setspace}
\onehalfspacing              % anderthalbzeilig (oder auch \doublespace)

%fancyBox
%\usepackage{fancybox}
\usepackage{placeins}

% Layout corrections (Schusterjungen)
\clubpenalty = 10000
% Layout corrections (Hurenkinder)
\widowpenalty = 10000
\displaywidowpenalty = 10000

% Figures
\usepackage{caption}
\usepackage[hypcap=true,labelformat=simple]{subcaption}
\renewcommand{\thesubfigure}{(\alph{subfigure})}

% Tables
\usepackage{booktabs}

% Frequently used column types
\newcolumntype{C}[1]{>{\centering\arraybackslash}p{#1}} % centering column type with fixed width
\newcolumntype{R}[1]{>{\raggedleft\arraybackslash}p{#1}} % right aligned column type with fixed width
\newcolumntype{L}[1]{>{\raggedright\arraybackslash}p{#1}} % left aligned column type with fixed width

% Shortcuts for referencing floats:
\newcommand{\fig}[1]{\figurename~\ref{#1}} %shortcut for a figure reference
\newcommand{\tab}[1]{Table~\ref{#1}} %shortcut for a table reference
\newcommand{\eq}[1]{(\ref{#1})} %shortcut for an equation reference
\newcommand{\lst}[1]{Listing~\ref{#1}} %shortcut for a listing reference
\newcommand{\sect}[1]{Section~\ref{#1}} %shortcut for a Section reference

\counterwithout{figure}{chapter}
\counterwithout{listing}{chapter}
\counterwithout{table}{chapter}

\begin{document}

\pagenumbering{roman}

\begin{titlepage}
	%select font without serifs
	\sffamily

	% Logo
	\begin{tabularx}{\textwidth}{@{}l@{}>{\raggedleft\arraybackslash}X@{}r@{}}
		\multirow{2}{*}{\includegraphics[width=6.8cm]{images/Logo_UniKassel}} &
		\raisebox{-1mm}{\small{Fachbereich Elektrotechnik/Informatik}}                                                       \\
		                                                                      & \raisebox{-1mm}{\small{\thesisdepartment}} &
	\end{tabularx}

	\vspace{2.5cm}

	\begin{center}
		% Title and subtitle
		\huge{\thesistitle}

		\vspace{3cm}

		\renewcommand{\baselinestretch}{1.3}
		\Large{\thesistype}

		\large
		\thesistypedesc
	\end{center}

	\vspace{1.5cm}
	\renewcommand{\baselinestretch}{1}
	\begin{table}[htpb]
		\centering
		\begin{tabular}{ll}
			\\
			Eingereicht von: & \thesisauthorname           \\
			Anschrift:       & \thesisauthorhomestreet     \\
			                 & \thesisauthorhometown       \\
			\\
			Matrikelnummer:  & \thesisauthormatrikelnumber \\
			E-Mail:          & \thesisauthoremail          \\
			\\
			Vorgelegt im:    & \thesisdepartment           \\
			\\
			Erstprüfer:      & \thesisfirstreviewer        \\
			Zweitprüfer:     & \thesissecondreviewer       \\
			\\
			Betreuer:        & \thesissupervisor           \\
			\\
			Eingereicht am:  & \thesisdate                 \\
		\end{tabular}
	\end{table}

	% font with serifs
	\rmfamily
\end{titlepage}

\setcounter{page}{0}
\chapter*{Eidesstattliche Erklärung}

Hiermit erkläre ich, dass ich die vorliegende Arbeit selbstständig und nur mit den nach der Prüfungsordnung der Universität Kassel zulässigen Hilfsmitteln angefertigt habe.
Die verwendete Literatur ist im Literaturverzeichnis angegeben.
Wörtlich oder sinngemäß übernommene Inhalte habe ich als solche kenntlich gemacht.

\vspace{1cm}

Kassel, \thesisdate

\begin{flushright}
  \underline{\hspace{7cm}} \\
  \thesisauthorname

  % \underline{\hspace{7cm}} \\
  % Weitere Autoren von Dokumentationen...
\end{flushright}

\chapter*{Zusammenfassung}

% Inhaltsverzeichnis und Kopfzeile
\addcontentsline{toc}{chapter}{Zusammenfassung}
\markboth{Zusammenfassung}{Zusammenfassung}

Diese Bachelorarbeit befasst sich mit der Erstellung von Konfigurationsdateien für die Open-Source-Software Kubernetes.
Dabei werden Probleme und Herausforderungen bei der Erstellung beleuchtet und bereits bestehende Lösungen
für die Probleme identifiziert. Dafür wird eine Stichwortsuche genutzt.
Im Rahmen der Suche wurde festgestellt, dass keine wissenschaftliche Arbeit existiert, die sich empirisch mit den Problemen und Herausforderungen
beim Erstellen der Konfigurationsdateien auseinandergesetzt.
Daher wurde eine Websuche durchgeführt, um diese zu erkennen.
Mit Hilfe eines Laborversuchs wurden Probleme erfasst, die von den Anwendungen nicht gelöst werden.
Hier konnte gezeigt werden, dass keine der existierenden Lösungen alle Probleme lösen kann.
Die Autovervollständigung und Validierung von Ressourcen auf dem Cluster in den Konfigurationsdateien wird von den Lösungen nicht angeboten.
Es wurde anschließend eine Anwendung implementiert, die diese Funktionalität für auswählte Cluster-Ressourcen bietet.
Die implementierte Anwendung konnte die funktionalen Anforderungen erfüllen.
Abschließend wurde diese mit den bestehenden Lösungen auf Basis von Kriterien zur Benutzerfreundlichkeit für statische Code-Analysetools verglichen.
Der Vergleich hat gezeigt, dass keine der untersuchten Anwendungen mehr als $50\%$ der Kriterien erfüllen konnte.
Die eigene Anwendung erfüllt nur 5 von 26 gestellten Kriterien.



\tableofcontents

\chapter*{Abkürzungsverzeichnis}

% Inhaltsverzeichnis und Kopfzeile
\addcontentsline{toc}{chapter}{Abkürzungsverzeichnis}
\markboth{Abkürzungsverzeichnis}{Abkürzungsverzeichnis}

% Auf diese Weise kann der Plural von unbekannten Wörtern definiert werden (für \acp{..})
\acrodefplural{ha}[HAs]{Hausaufgaben}
\acrodefplural{rq}[RQs]{Forschungsfragen}

\begin{acronym}[XXXXXX] % Anzahl der X gibt an, welche Breite das längste Acronym hat
    % Es werden nur Akronyme übernommen, die auch verwendet werden
    \acro{api}[API]{Application Programming Interface}
    \acro{ast}[AST]{Abstract Syntax Tree}
    \acro{css}[CSS]{Cascading Style Sheet}
    \acro{gui}[GUI]{Graphical User Interface}
    \acro{ha}[HA]{Hausaufgabe}
    \acro{html}[HTML]{Hypertext Markup Language}
    \acro{http}[HTTP]{Hypertext Transfer Protocol}
    \acro{https}[HTTPS]{Hypertext Transfer Protocol Secure}
    \acro{id}[ID]{Identifier}
    \acro{ide}[IDE]{Integrated Development Environment}
    \acro{ip}[IP]{Internet Protocol}
    \acro{json}[JSON]{JavaScript Object Notation}
    \acro{jvm}[JVM]{Java Virtual Machine}
    \acro{mime}[MIME]{Multipurpose Internet Mail Extensions}
    \acro{npm}[npm]{Node Package Manager}
    \acro{rest}[REST]{Representational State Transfer}
    \acro{rpc}[RPC]{Remote Procedure Call}
    \acro{rq}[RQ]{Forschungsfrage}
    \acro{sdk}[SDK]{Software Development Kit}
    \acro{ssh}[SSH]{Secure Shell}
    \acro{ui}[UI]{User Interface}
    \acro{url}[URL]{Uniform Resource Locator}
    \acro{xml}[XML]{Extensible Markup Language}
    \acro{lsp}[LSP]{Language Server Protokoll}
    \acro{vscode}[VS Code]{Visual Studio Code}
    \acro{vsce}[vsce]{Visual Studio Code Extensions}
    \acro{yaml}[YAML]{Yet Another Markup Language}

    \vspace{\parskip}

    \acro{bzw}[bzw.]{beziehungsweise}
    \acro{dh}[d.h.]{das heißt}
    \acro{etc}[etc.]{et cetera}
    \acro{idR}[i.d.R.]{in der Regel}
    \acro{oBdA}[o.B.d.A.]{ohne Beschränkung der Allgemeinheit}
    \acro{sic}[sic]{sic erat scriptum}
    \acro{usw}[usw.]{und so weiter}
    \acro{uU}[u.U.]{unter Umständen}
    \acro{vgl}[vgl.]{vergleiche}
    \acro{zB}[z.B.]{zum Beispiel}
\end{acronym}


\pagebreak
\pagenumbering{arabic}

% Hier weitere Kapitel einfügen
\chapter{Einleitung}\label{ch:introduction}

Die Skalierbarkeit und Ausfallsicherheit von IT-Systemen wird mit der zunehmenden Digitalisierung von Prozessen immer wichtiger.
Dies zeigt auch das Problem bei Erreichen des KfW-Portals am 26.09.23 für das Förderprogramm ``Solarstrom für Elektroautos''~\cite{kfw-not-reachable}.
Skalierbarkeit bezeichnet die Fähigkeit eines Systems bei erhöhter Nachfrage Ressourcen hinzuzufügen, um die gestiegende Nachfrage zu bewältigen und
diese Ressourcen bei sinkender Nachfrage wieder freizugeben. Das System kann seine Kapizität an wechselnde Anforderungen anpassen~\cite{it-system-scaling}.
Um die Skalierbarkeit von IT-Systemen zu realisieren, bietet sich die Software Kubernetes an.
Bei Kubernetes handelt es sich um eine Software zur Bereitstellung, Verwaltung und Skalierung von Anwendungen in Containern~\cite{kubernetes-overview}.
Ein Kubernetes Cluster besteht aus sogenannten Kubernetes Objekten, die den aktuellen Zustand des Clusters beschreiben. Dabei geben die Kubernetes Objekte an,
welche Anwendungen auf dem Cluster laufen, welche Ressourcen diesen zur Verfügung stehen und wie sich die Anwendungen, \acs{zB} im Fehlerfall, verhalten~\cite{kubernetes-objects}.

\section{Problemstellung}\label{sec:problem}

Die Definition von Kubernetes Objekten kann in deklarativen Konfigurationsdateien im \acs{yaml}-Format vorgenommen werden~\cite{kubernetes-config-declarative}.
Das Schreiben und Pflegen dieser Konfigurationsdateien ist schwierig und fehleranfällig, insbesondere wenn die Größe und Komplexität des Clusters zunimmt~\cite{kubernetes-config-problems} TODO: find more sources.
Im Rahmen dieser Arbeit werden die Möglichkeiten der Autovervollständigung und Validierung von Kubernetes Konfigurationsdateien untersucht.
Dabei werden folgende Fragen beantwortet:
\begin{enumerate}
    \item Wie kann die Autovervollständigung und Validierung von Kubernetes Konfigurationsdateien in eine \acs{ide} integriert werden?
    \item Welche Vor- und Nachteile hat die Implementierung im Vergleich zu bereits bestehenden Lösungen?
\end{enumerate}

\section{Methodik}
Um die in Kapitel~\ref{sec:problem} gestellten Fragen zu beantworten, wird eine \acs{ide}-Erweiterung entwickelt, welche die Validierung und Autovervollständigung von
Kubernetes Konfigurationsdateien ermöglicht. Anschließend wird die entwickelte Lösung mit bereits bestehenden Lösungen verglichen.
Hierbei werden ausgwählte Kriterien zur Beurteilung aus~\cite[A Large-Scale Study of Usability Criteria Addressed by Static Analysis Tools]{usability-criteria-static-analysis-tools} herangezogen.
Außerdem wird die Autovervollständigung ausgewählter Lösungen in einem Laborversuch untersucht.
Die Ergebnisse der Untersuchungen werden abschließend evaluiert.

\section{Aufbau der Arbeit}
Im folgenden Kapitel werden die technischen Grundlagen beschrieben. In Kapitel~\ref{ch:implementation} werden die Schritte zur Implementierung erläutert.
Der Vergleich der Implementierung mit ähnlichen Lösungen, sowie der Laborversuch erfolgen in Kapitel~\ref{ch:comparison}.
Die Vor- und Nachteile werden der verschiedenen Lösungen im anschließenden Kapitel evaluiert.
Abschließend folgt das Fazit der Arbeit und ein Ausblick auf weitere mögliche Schritte.
\chapter{Technische Grundlagen}\label{ch:technical-basics}

\section{Virtualisierung}

Mit Hilfe von Virtualisierung lassen sich mehrere simulierte Umgebungen aus einem einzigen Hardwaresystem, im folgenden \textit{Host} genannt, erstellen.
Dazu wird Software genutzt, die die physische Hardware emuliert~\cite{virtualization-techtarget}. Diese Software wird auch \textit{Hypervisor} genannt.
Der \textit{Hypervisor} abstrahiert die Hardware Ressourcen des \textit{Hosts} und stellt diese als einen gemeinsamen Pool zur Verfügung.
Die physische Hardware des \textit{Hosts} kann somit in mehrere isolierte Umgebungen, auch virtuelle Maschinen genannt, aufgeteilt werden.
Mithilfe des \textit{Hypervisors} können mehrere unterschiedliche Betriebssysteme parallel ausgeführt werden~\cite{virtualization-red-hat, hypervisor-red-hat}.
Das heißt, dass jede virtuelle Maschine ein eigenes Betriebssystem benötigt, um eine Anwendung auszuführen.

Sollen zwei verschiedene Anwendungen isoliert voneinander ausgeführt werden, werden dafür zwei virtuelle Maschinen mit jeweils einem eigenen Betriebssystem benötigt.
Jede virtuelle Maschine führt dabei ihren eigenen Satz von Systemprozessen aus, was neben der Ausführung der Anwendung zusätzliche Ressourcen erfordert~\cite{kubernetes-in-action-introduction}.
Um dieses Problem zu lösen wurde die Containervirtualisierung eingeführt.

\begin{figure}
      \centering
      \begin{subfigure}{.5\textwidth}
            \centering
            \includesvg{images/full-virtualization.svg}
            \caption{Vollvirtualisierung}
            \label{fig:full-virtualization}
      \end{subfigure}%
      \begin{subfigure}{.5\textwidth}
            \centering
            \includesvg{images/container-virtualization.svg}
            \caption{Containervirtualisierung}
            \label{fig:container-virtualization}
      \end{subfigure}
      \caption{Vergleich Voll- und Containervirtualisierung}
      \label{fig:virtualization-comparison}
\end{figure}

Abbildung~\ref{fig:virtualization-comparison} zeigt einen Vergleich zwischen der soeben eingeführten Virtualisierung mittels \textit{Hypervisor} und
der Containervirtualisierung.
Abbildung~\ref{fig:full-virtualization} zeigt die Komponenten einer Vollvirtualisierung, bei der der \textit{Hypervisor}
eine auf dem \textit{Host}-Betriebssystem installierte Software ist.
\\
Die Komponenten einer Containervirtualisierung sind auf Abbildung~\ref{fig:container-virtualization} dargestellt.
Im Gegensatz zur Vollvirtualisierung können Anwendungen über sogennante \textit{Container} voneinander getrennt werden.
Ein \textit{Container} ist ein einziger, isolierter Prozess, der im Betriebssystem des Hosts läuft und nur die Ressourcen nutzt, die von der Anwendung benötigt werden.
Es ensteht so kein Mehraufwand für Systemprozesse, wie bei der Vollvirtualisierung~\cite{kubernetes-in-action-introduction}.
Um die Bereitstellung von \textit{Containern} zu ermöglichen wird eine \textit{Container Runtime} benötigt~\cite{kubernetes-overview}.
\\
Eine \textit{Container Runtime} ist eine Softwarekomponente, die die Ausführung von \textit{Containern} auf einem Host-Betriebssystem ermöglicht.
Die Aufgaben einer \textit{Container Runtime} sind dabei das Erstellen, Starten, Stoppen und Löschen von \textit{Containern},
das Überwachen der Systemressourcen des Hosts, das Isolieren von Systemressourcen für einen \textit{Container}
und das Laden von \textit{Container Images} aus einem \textit{Respository}~\cite{container-runtime-aquasec}.
\\
Ein \textit{Container Image} enthält Software und alle Ressourcen, die die Software zur Ausführung benötigt.
Neben dem Programmcode und den Bibliotheken, die die Anwendung zur Laufzeit nutzt, können dies auch weitere Programme sein, auf die zugegriffen werden muss.
\textit{Container Images} werden über ein zentrales \textit{Respository} verwaltet und können darüber bei Bedarf heruntergeladen werden~\cite{container-image-dev-insider}.

\section{Kubernetes}

Kubernetes (K8s) wurde ursprünglich von Google entwickelt und 2014 als Open Source Software veröffentlicht.
Kubernetes ist eine Software zur Bereitstellung, Verwaltung und Skalierung von Anwendungen in Containern. Der Name kommt aus dem Griechischen und bedeutet Steuermann.
Die Abkürzung K8s setzt sich aus dem ersten und letzten Buchstaben von Kubernetes, sowie der Anzahl der Buchstaben die diese voneinander trennen, zusammen~\cite{kubernetes-google-cloud}.

\subsection{Architektur}

Kubernetes abstrahiert die zur Verfügung stehende Hardware und fasst mehrere physische oder virtuelle Maschinen,
die als \textit{Nodes} bezeichnet werden, zu einem Cluster zusammen.
Die zugrundeliegende Infrastruktur wird dabei ausgeblendet und Kubernetes betrachtet die \textit{Nodes} als eine einzige, große Ressource.
\\
Jedes Cluster besteht aus verschiedenen Typen von \textit{Nodes}.
Für ein funktionierendes Kubernetes Cluster wird mindestens ein \textit{Master Node (Control Plane)} und ein \textit{Worker Node} benötigt~\cite{kubernetes-in-action-introduction}.
Auf Abbildung~\ref{fig:kubernetes-architecture} ist der minimale Aufbau eines Kubernetes-Clusters schematisch dargestellt. Es folgt eine
Beschreibung der Aufgaben des \textit{Master} und \textit{Worker Nodes}, sowie eine Beschreibung der einzelnen Komponenten.

\begin{figure}[htp] % htp = hier (h), top (t), oder auf einer eigenen Seite (p).
      \centering
      \includesvg[width=0.5\textwidth]{images/kubernetes-architecture.svg} % width immer angeben!
      \caption{Komponenten einer Kubernetes-Clusters}
      \label{fig:kubernetes-architecture}
\end{figure}

Das \textit{Master Node} ist für die Administration des Clusters zuständig. Es koordiniert und überwacht die \textit{Worker Nodes} des Clusters.
Seine Aufgaben sind unter anderem die Planung und Skalierung von Anwendungen und die Aufrechterhaltung des Zustands des Clusters~\cite{kubernetes-cluster-vmware}.
Das \textit{Master Node} besteht aus den vier folgenden Komponenten:
\begin{description}
      \item[API-Server]
            Der \textit{API-Server} koordiniert alle Prozesse zwischen dem \textit{Master Node} und den Komponenten der \textit{Worker Nodes}.
            Er stellt außerdem die \textit{Kubernetes-API}, welche als Zugang zum Cluster von außen dient, bereit und authentifiziert und authorisiert die Zugriffe.
            Kommandozeilen-Tools wie kubectl oder Drittanbieter-Software können so auf das Cluster zugreifen~\cite{kubernetes-devopscube}.
      \item[Scheduler]
            Der \textit{Scheduler} ist für die Verteilung von \textit{Containern} auf die \textit{Worker Nodes} zuständig.
            Er berücksichtigt dabei die zur Verfügung stehenden Ressourcen der einzelnen \textit{Nodes}~\cite{kubernetes-cluster-ionos}.
      \item[Controller-Manager]
            Der Controller-Manager handhabt clusterweite Funktionen zur Aufrechterhaltung des Zustands des Clusters. Er behandelt \ac{zB} den Ausfall von \textit{Nodes}
            im Cluster~\cite{kubernetes-in-action-introduction}.
      \item[etcd]
            Die \textit{etcd} ist eine konsistenter Key-Value Datenbank, die in verteilten Systmen eingesetzt werden kann~\cite{etcd-ibm}.
            Bei Kubernetes speichert \textit{etcd} die gesamte Clusterkonfiguration~\cite{kubernetes-in-action-introduction}.
\end{description}

Ein \textit{Worker Node} ist für die Ausführung von den Anwendungen in \textit{Containern} zuständig.
Es führt Aufgaben aus, die diesem von dem \textit{Master Node} zugewiesen werden~\cite{kubernetes-cluster-vmware}.
Ein \textit{Worker Node} besteht aus den folgenden zwei Komponenten:
\begin{description}
      \item[Kubelet]
            Das \textit{Kubelet} steuert die \textit{Container} auf dem \textit{Worker Node}~\cite{kubernetes-in-action-introduction}.
            Zur Ausführung der \textit{Container} wird eine \textit{Container-Runtime} genutzt.
            Das \textit{Kubelet} steht im Austausch mit dem \textit{API-Server} des \textit{Master Nodes}. Es überwacht \ac{zB} die laufenden Container
            und meldet deren Status, Ereignisse und Ressourcenverbrauch an den \textit{API-Server}~\cite{kubernetes-in-action-chapter-11}.
      \item[kube-proxy]
            Der \textit{kube-proxy} ist für die Netzwerkkonnektivtät des \textit{Worker Nodes} verantwortlich.
            Er ist für die Umsetzung von Netzwerkregeln und die Lastenverteilung zuständig~\cite{kubernetes-in-action-introduction,kubernetes-cluster-vmware}.
\end{description}

\subsection{Objekte in Kubernetes}

Kubernetes Objekte sind persistente Entitäten eines Kubernetes Systems. Sie repräsentieren den aktuellen Zustand des Clusters.
Die Kubernetes Objekte können beschreiben, welche Anwendungen auf dem Cluster in \textit{Containern} laufen,
welche Ressourcen den Anwendungungen zur Verfügung stehen und
wie sich die Anwendungen verhalten, \ac{zB} unter welchen Umständen eine Anwendung neu gestartet werden soll~\cite{kubernetes-objects}.
\\
Eines der Kernkonzepte von Kubernetes ist der sogennante \textit{Desired State}~\cite{kubernetes-objects-kodekloud}.
Der Nutzer definiert mit Hilfe Kubernetes Objekte den angestrebten Zustand (\textit{Desired State}) des Clusters und Kubernetes
sorgt für die Herstellung und Aufrechterhaltung dieses Zustandes.
\\
Die Definition eines Kubernetes Objektes erfolgt in Form einer \textit{Manifest}-Datei im \ac{yaml}-Format.
In jeder \textit{Manifest}-Datei müssen folgende Attribute definiert sein~\cite{kubernetes-objects}:

\begin{itemize}
      \setlength\itemsep{-0.5cm}
      \item
            \textbf{apiVersion} Gibt an welche Version der \textit{Kubernetes API} genutzt wird, um das Objekt zu erstellen.
      \item
            \textbf{kind}  Gibt die Art des Kubernetes Objektes an.
      \item
            \textbf{metadata}  Enhält Daten die das Objekt eindeutig identifizieren, \ac{zB} ein Name und eine ID\@.
      \item
            \textbf{spec} Beschreibt den \textit{Desired State} des Objektes.
\end{itemize}

Kapitel~\ref{ch:implementation} und Kapitel~\ref{ch:comparison} nutzen spezifische Kubernetes Objekte und Konzepte die im Folgenden
zum besseren Verständnis beschrieben werden:

\begin{description}
      \item[Pod]
            Ein \textit{Pod} ist eine Gruppe von \textit{Containern}, die sich auf dem selben \textit{Worker Node} befinden.
            Es wird sichergestellt, dass sich alle \textit{Container}, die in einem Pod zusammengefasst,
            immer auf dem selben \textit{Worker Node} ausgeführt werden~\cite{kubernetes-objects-kodekloud,kubernetes-in-action-chapter-3}.
      \item[Replica Set]
            Ein \textit{Replica Set} führt eine spezifizierte Anzahl von \textit{Pods} aus und sorgt dafür, dass diese Anzahl konstant bleibt.
            Sollte ein \textit{Pod} ausfallen, sorgt das \textit{Replica Set} dafür, dass ein neuer gestartet wird. Bei hoher Ressourceauslastung
            nutzt das \textit{Replica Set} eine Lastenverteilung, um alle \textit{Pods} gleichmäßig auszulasten~\cite{kubernetes-replica-set-kubematic}.
      \item[Deployment]
            Mit Hilfe eines \textit{Deployments} ist es möglich die Anwendungen in \textit{Pods}, die einem \textit{Replica Set} zugeordnet sind, zu aktualisieren.
            In der \textit{Manifest}-Datei kann dazu der \textit{Desired State} definiert werden.
            Das Ausrollen der Aktualisierung erfolgt für alle \textit{Pods} schrittweise.
            Tritt bei der Aktualisierung eines \textit{Pods} ein Fehler auf, so kann dieser neu erstellt oder wenn nötig auf den alten Stand zurückgesetzt werden.
            Kubernetes sorgt dafür, dass die Ausfallzeit der Anwendung so gering wie möglich ist~\cite{kubernetes-in-action-chapter-9,kubernetes-deployment-vmware}.
      \item[Secret]
            Bei einem \textit{Secret} handelt es sich um ein Kubernetes Objekt, das sensible Daten, wie \ac{zB} Passwörter oder SSH-Schlüssel, verschlüsselt speichert.
            Sensible Daten müssen so nicht in der \textit{Manifest}-Datei des \textit{Pods} und im Programmcode der Anwendung gespeichert werden~\cite{kubernetes-secrets,kubernetes-secrets-medium}.
      \item[Namespace]
            Ein \textit{Namespace} ermöglicht dem Nutzer ein physisches Kubernetes-Cluster in verschiedene virtuelle Cluster aufzuteilen.
            Sie können zum Einsatz kommen, wenn sich verschiedene Teams ein Kubernetes-Cluster teilen~\cite{kubernetes-namespaces-vmware}.
      \item[Label]
            Ein \textit{Label} ist ein Schlüssel-Wert-Paar, das einem Kubernetes Objekt, wie \ac{zB} einem \textit{Pod}, zugeordnet werden kann.
            \textit{Labels} können im \textit{metadata} Attribut der \textit{Manifest}-Datei eines Objektes angegeben werden.
            Der Schlüssel eines \textit{Labels} muss für ein Objekt eindeutig sein. Ein \textit{Label} kann mehreren Objekten zugeordnet werden.
            Der Nutzer hat die Möglichkeit eine Menge von Objekten über ein \textit{Label} zu selektieren~\cite{kubernetes-labels}.
\end{description}

\section{\acs{lsp}}

Das \ac{lsp} wurde 2015 von Microsoft für den Editor Visual Studio Code entwickelt.
Bei einem Editor handelt es sich in diesem Kontext um ein Computerprogramm zur Erstellung und Bearbeitung von Textdateien~\cite{editor-definition}.
Das Ziel war es, eine Schnittstelle bereitzustellen, die unter anderem die Autovervollständigung, die Code Navigation und das Syntaxhighlighting
für jeden Editor vereinheitlicht~\cite{lsp-witekio}.
Die Fähigkeit eines Editors Autovervollständigung, Code Navigation und Syntaxhighlighting zu ermöglichen wird im Folgenden als Sprachunterstützung bezeichnet.
Das \ac{lsp}-Protokoll soll vermeiden, dass eine Implementierung der Sprachunterstützung für jeden Editor seperat durchgeführt werden muss.
Abbildung~\ref{fig:without-lsp} veranschaulicht dieses Problem.

\begin{figure}[htp] % htp = hier (h), top (t), oder auf einer eigenen Seite (p).
      \centering
      \includesvg[width=0.5\textwidth]{images/without-LSP.svg} % width immer angeben!
      \caption{Sprachunterstützung der Editoren ohne das \acs{lsp}}
      \label{fig:without-lsp}
\end{figure}

Als Lösung für das oben stehende Problem definiert das \ac{lsp} einen \ac{lsp}-Client und einen \ac{lsp}-Server.
Der \ac{lsp}-Client ist eine Erweiterung des Editors und sendet \ac{zB} Änderungen an Dateien zum Server.
Der \ac{lsp}-Server enthält auf der anderen Seite die Logiken zur Sprachunterstützung~\cite{lsp-witekio}.
Abbildung~\ref{fig:with-lsp} zeigt, dass mit Hilfe des \ac{lsp} die Sprachunterstützung nur einmal auf Server Seite entwickelt wird.
Für die verschiedenen Editoren muss folglich nur eine Client Erweiterung implementiert werden, um alle Sprachen zu unterstützen.

\begin{figure}[htp] % htp = hier (h), top (t), oder auf einer eigenen Seite (p).
      \centering
      \includesvg[width=0.5\textwidth]{images/with-lsp.svg} % width immer angeben!
      \caption{Sprachunterstützung der Editoren mit \acs{lsp}}
      \label{fig:with-lsp}
\end{figure}

Die Kommunikation zwischen \ac{lsp}-Client und \ac{lsp}-Server erfolgt über Nachrichten, die jeweils aus einem \textit{Header}- und aus einem \textit{Content}-Teil bestehen.
Der \textit{Header}-Teil enthält zum einen die Länge des \textit{Content}-Teils in Byte. Zum Anderen enthält dieser eine Typenbezeichnung,
die den \ac{mime}-Typ des \textit{Content}-Teils angibt. Der \textit{Content}-Teil enhält den Inhalt der Nachricht. Die Struktur des Inhalts wird durch
das \acs{json}-\acs{rpc} Protokoll vorgegeben~\cite{lsp-microsoft,lsp-medium}.
\\
Bei \ac{rpc} handelt es sich um ein Programmierparadigma, dass die Entwicklung von verteilten Anwendungen vereinfachen soll.
Das Konzept ist dabei, dass ein Programm auf einem Client eine Methode auf eineem entfernten System, einem Server, aufrufen kann,
ohne die zugrundliegenden Netzwerkdetails zu kennen \cite{rpc-computer-weekly}.

Das \acs{json}-\acs{rpc} Protokoll definiert für die Kommunikation drei Arten von Nachrichten~\cite{lsp-microsoft}:
\begin{description}
      \item[Request Message] Die \textit{Request Message} ist eine Nachricht vom Client an den Server.
            Sie enthält eine ID zur eindeutigen Identifikation der Nachricht und die Methode, die auf dem Server aufgerufen werden soll, sowie die Methodenparameter.
            Der Server muss eine Antwort an den Client zurücksenden.
      \item[Response Message] Die \textit{Response Message} ist die Antwort des Servers an den Client. Sie enthält die ID der dazugehörigen \textit{Request Message} und das Ergebnis der aufgerufenen Methode.
            Im Fehlerfall enthält die \textit{Response Message} kein Ergebnis, sondern eine Beschreibung des Fehlers.
      \item[Notification Message]  Die \textit{Notification Message} kann vom Client oder vom Server versendet werden.
            Der Empfänger muss auf diese Art von Nachricht nicht antworten.
            Sie enhält den Ereignis-Typ, sowie die Ereignis-Parameter.
\end{description}

Die Methoden und Ereignisse die Client und Server zur Verfügung stellen werden als \textit{Capabilities} (Fähigkeiten) bezeichnet~\cite{lsp-medium}.
Das \ac{lsp} gibt \textit{Capabilities} vor, die jeder Client implementieren muss:
\begin{description}
      \item[textDocument/didOpen] Eine \textit{Notification Message} die dem Server signalisiert, dass ein neues Dokument geöffnet wurde.
      \item[textDocument/didChange] Eine \textit{Notification Message} die dem Server eine Änderung in einem Dokument anzeigt.
      \item[textDocument/didClose]  Eine \textit{Notification Message} die dem Server anzeigt, dass eine Dokument geschlossen wurde.
\end{description}

Beim Initialisieren der Verbindung zwischen Client und Server müssen diese jeweils der anderen Stelle ihre \textit{Capabilities} mitteilen.
Der Client sendet eine Nachricht zur Intialisierung mit seinen \textit{Capabilities} an den Server.
Dieser antwortet wiederum mit einer Nachricht, die dessen \textit{Capabilities} enthält.
Nach Erhalt dieser Nachricht sendet der Client erneut eine Nachricht an den Server um zu signalisieren, dass die Intialisierung abgeschlossen ist.
Anschließend kann der Nachrichtenaustausch zwischen Client und Server stattfinden~\cite{lsp-microsoft,lsp-medium}.

\begin{figure}[htp] % htp = hier (h), top (t), oder auf einer eigenen Seite (p).
      \centering
      \includesvg[width=0.5\textwidth]{images/lsp-sequence-diagram.svg} % width immer angeben!
      \caption{Sequenzdiagramm einer \acs{lsp}-Kommunikation}
      \label{fig:lsp-sequence-diagram}
\end{figure}

Abbildung~\ref{fig:lsp-sequence-diagram} zeigt ein Sequenzdiagramm einer Beispielkommunikation zwischen einem \ac{lsp}-Client und einem \ac{lsp}-Server.
Dabei erfolgt die Intialisierung der Verbindung mit Mitteilung der \textit{Capabilities}.
Danach teilt der Client dem Server in Form einer \textit{Notification Message} mit, dass ein Dokument geöffnet wurde.
Zuletzt sendet der Server dem Client eine \textit{Notification Message} mit dem Ergebnis einer Validierung, das \acs{zB} die Syntax Fehler in einer Datei enhalten könnte.

\section{\acs{json}-Schema}

\acs{json}-Schema ist eine deklarative Sprache, die die Validierung des Formates und der Struktur eines \ac{json}-Objektes ermöglicht.
Mit Hilfe des \acs{json}-Schema können Regeln und Einschränkungen definiert werden, um den Aufbau eines validen \ac{json}-Objektes zu beschreiben~\cite{json-schema-description}.

\begin{listing}[htp]
      \begin{minted}{json}
      {
            "product": {
                  "name": "Widget",
                  "price": 10.99
            }
      }
      \end{minted}
      \caption{Ein \acs{json}-Objekt}
      \label{lst:json-object}
\end{listing}

\begin{listing}[h]
      \begin{minted}{json}
      {
            "$schema": "https://json-schema.org/draft/2020-12/schema",
            "title": "Product",
            "type": "object",
            "properties": {
                "product": {
                    "type": "object",
                    "properties": {
                        "name": {
                            "type": "string"
                        },
                        "price": {
                            "type": "number", 
                            "minimum": 0
                        },
                    },
                    "required": ["name", "price"]
                }
            },
            "required": ["product"]
      }
      \end{minted}
      \caption{Ein \acs{json}-Schema}
      \label{lst:json-schema}
\end{listing}

Die Regeln und Einschränkungen werden mit einem Satz von Schlüsselwörtern definiert und Objekten im \acs{json}-Format definiert.
\\
Auf Listing~\ref{lst:json-object} ist ein Beispiel eines \acs{json}-Objektes abgebildet, dass mit Hilfe des auf Listing~\ref{lst:json-schema}
dargestellten \acs{json}-Schemas validiert werden kann.
\\
Es folgt eine Erklärung der verwendeten Schlüsselwörter~\cite{json-schema-creation}:

\begin{itemize}
      \setlength\itemsep{-0.5cm}
      \item
            \textbf{\$schema} Gibt die Version der \acs{json}-Schema Spezifikation. Unterschiedliche Versionen der Spezifikation
            unterstützen andere Schlüsselwörter.
      \item
            \textbf{type} Legt den Typ des Attributes im \acs{json}-Format fest.
      \item
            \textbf{properties} Gibt an, welche Attribute ein \acs{json}-Objekt enthalten kann.
      \item
            \textbf{required} Legt fest, welche Attribute ein \acs{json}-Objekt enthalten muss.
      \item
            \textbf{minimum} Gibt den kleinsten Wert an, den ein numerischer Wert enthalten darf.
\end{itemize}

Anhand der im \acs{json}-Schema enthaltenen Informationen kann ein Algorithmus die Validierung des Objektes durchführen.
Die \acs{json}-Schema Spezifikation enthält noch weitere Schlüsselwörter, die aufgrund des Umfangs nicht erläutert werden.
\chapter{Implementierung}\label{ch:implementation}

Um die in Kapitel~\ref{ch:introduction} gestellten Forschungsfragen zu beantworten, wurde ein \ac{lsp}-Server und ein \ac{lsp}-Client entwickelt,
welche die Autovervollständigung und Validierung von Kubernetes \textit{Manifest}-Dateien ermöglichen.
Die beiden Komponenten wurden als installierbare Erweiterung für den Quellcode-Editor \ac{vscode} bereitgestellt.

\section{Technologien}

\begin{description}
      \item[NodeJS]
            Bei \textit{NodeJS} handelt es sich um eine plattformübergreifende Laufzeitumgebung, mit welcher Javascript-Programme entwickelt werden können,
            die unabhängig von Host-Anwendungen, wie \ac{zB} dem Webbrowser laufen.
            Das Konzept des ``non-blocking I/O'' wird von \textit{NodeJS} für die parallele und asynchrone Ausführung von Code verwendet.
            Programme müssen dabei nicht auf die Ergebnisse einer Ein-/Ausgabe-Funktion warten,
            sondern werden über ein Ereignis über die Beendigung informiert~\cite{node-js-dev-insider,node-js-about}.
            Im Rahmen des Projektes wurden mit Hilfe von NodeJS der \ac{lsp}-Client und \ac{lsp}-Server implementiert.
      \item[Typescript]
            \textit{Typescript} ist eine Skriptsprache, die eine syntaktische Obermenge von Javascript ist. Das heißt, dass \textit{Typescript} und Javascript
            die gleiche Syntax besitzen, Typescript jedoch noch über zusätzliche Syntax verfügt.
            \textit{Typescript} erlaubt dem Programmierer \ac{zB} die statische Typisierung von Variablen, Funktionsparametern und Rückgabetypen.
            Die Code-Qualität kann dadurch erhöht werden. Ein Compiler übersetzt \textit{Typescript}-Code in lauffähigen Javascript-Code, welcher dann von einer
            wie \ac{zB} \textit{NodeJS} ausgeführt werden kann~\cite{typescript-kinsta,typescript-doubleslash-blog}.
            Der \ac{lsp}-Client und \ac{lsp}-Server wurden mit der Skriptsprache \textit{Typescript} entwickelt.
      \item[Docker for Desktop]
            Mit \textit{Docker for Desktop} lassen sich Anwendungen in \textit{Containern} zusammenstellen, teilen und ausführen.
            Das Programm ist für Mac, Linux und Windows Umgebungen verfügbar und bietet ein \ac{gui} zur Verwaltung der \textit{Container} an~\cite{docker-for-desktop-overview}.
            \textit{Docker for Desktop} enthält außerdem ein Kubernetes-Cluster, welches lokal auf einem einzigen \textit{Node} ausgeführt wird~\cite{docker-for-desktop-kubernetes}.
            Die Kubernetes-API des Clusters wurde bei der Implementierung des \ac{lsp}-Server verwendet, um
            die bestehenden Objekte auf dem Cluster abzufragen und diese für die Autovervollständigung zu verwenden.
      \item[kubectl]
            \textit{kubectl} ist ein Kommandozeilentool, welches zum erstellen, untersuchen, aktualisieren und löschen von Kubernetes Objekten, sowie zum Ausrollen
            von Anwendungen und verwalten der Cluser-Ressourcen verwendet werden kann. Es kommuniziert dazu mit der Kubernetes-API.\@
            Bei der Entwicklung wurde \textit{kubectl} verwendet um den momentanen Zustand des Clusters abzufragen,
            um zu verifizieren, welche Objekte der \ac{lsp}-Server zurückgeben muss.
            Außerdem konnte der ``verbose''-Parameter der \textit{kubectl} verwendet, um die benötigten Anfragen an die Kubernetes-API anzeigen zu lassen.
            Diese wurden für die Implementierung des \ac{lsp}-Servers genutzt.
\end{description}

\section{Vorgehen}

\section{Anforderungen}

\section{Bereitstellung des Clusters}

\section{Erstellen der \acs{json}-Schema Dateien aus der OpenAPI-Spezifikation von Kubernetes}

\section{Umwandlung des \acs{ast}}

\section{Autovervollständigung mit Hilfe des \acs{json}-Schemas}

\section{Autovervollständigung mit Hilfe der Kubernetes-\acs{api}}


\chapter{Vergleich}\label{ch:comparison}

\chapter{Evaluation}\label{ch:evaluation}

\section{Bewertung anhand der Anforderungen}

\section{Vergleich mit anderen Tools anhand von gegebenen Kriterien}
\chapter{Fazit}\label{ch:conclusion}


% Die nächsten zwei Zeilen sind optional, sie sorgen dafür dass alles nach dem Inhalt wieder mit römischen Zahlen nummeriert wird.
\pagenumbering{roman}
\addtocounter{page}{4} % Dies ist die Anzahl der Seiten vor der Einleitung, muss möglicherweise angepasst werden, wenn das Inhaltsverzeichnis mehrere Seiten umfasst.

\bibliographystyle{unsrt}
\bibliography{main}

\listoffigures
\listoftables
\renewcommand{\listoflistingscaption}{Listing-Verzeichnis}
\listoflistings

\appendix
\chapter{Anhang}\label{ch:appendix}

\begin{figure}[htp] % htp = hier (h), top (t), oder auf einer eigenen Seite (p).
    \centering
    \includesvg[width=0.8\textwidth]{images/language-service-interface.svg} % width immer angeben!
    \caption{Klassendiagramm des Interface ``LanguageService''}
    \label{fig:language-service-interface-defition}
\end{figure}

\vspace{0.5cm}

\begin{figure}[htp] % htp = hier (h), top (t), oder auf einer eigenen Seite (p).
    \centering
    \includesvg[width=0.35\textwidth]{images/class-diagram-json-document.svg} % width immer angeben!
    \caption{Klassendiagramm der Klasse ``JSONDocument''}
    \label{fig:class-diagram-json-document}
\end{figure}

\begin{figure}[htp] % htp = hier (h), top (t), oder auf einer eigenen Seite (p).
    \centering
    \includesvg[angle=90, width=0.73\paperheight]{images/class-diagram-lsp-server.svg} % width immer angeben!
    \caption{Klassendiagramm \acs{lsp}-Server}
    \label{fig:class-diagram-lsp-Server}
\end{figure}

\begin{figure}[htp] % htp = hier (h), top (t), oder auf einer eigenen Seite (p).
    \centering
    \includesvg[width=0.8\textwidth]{images/flowchart-get-node-from-offset.svg} % width immer angeben!
    \caption{Ablaufdiagramm der Methode ``getNodeFromOffset''}
    \floatfoot{Im Ablaufdiagramm steht ``CP'' für die Position des Cursors und ``K'' kann einen Knoten aus dem \ac{ast} eines
        \ac{yaml}-\textit{Documents} enthalten.}
    \label{fig:flowchart-get-node-from-offset}
\end{figure}






\end{document}
